\section{Positioning relative to DeSa2VA and parametric lanes}
\label{app:position}

Table~\ref{tab:pos} summarizes what we take from the closest architectural and geometric lines and what we do not claim.
DeSa2VA is not re-run on FIRC2234: it is an architectural fork of Sa2VA trained on RefCOCO-family and video-RES mixtures, and PolyLaneNet / B\'ezierLaneNet are lane detectors, not MLLM referring models.
The table is a mechanistic comparison, not a leaderboard.

\begin{table*}[t]
\centering
\footnotesize
\setlength{\tabcolsep}{4pt}
\caption{Mechanistic positioning. Numbers in this paper remain FIRC-only; DeSa2VA and lane-detector numbers stay on their native benchmarks.}
\label{tab:pos}
\begin{tabularx}{\textwidth}{@{}lXXX@{}}
\toprule
 & DeSa2VA \citep{dang2025desa2va} & Poly/B\'ezier lane heads \citep{tabelini2021polylanenet,feng2022bezierlanenet} & This paper \\
\midrule
Special token & shared \Seg{}, then split & none (CNN detector) & shared \Seg{} (undecoupled) \\
Mask carrier & SAM2 after decoupling & polynomial / B\'ezier curve & $\phih \to \cR$ (SAM2 bypassed) \\
Reported failure & extra text into entangled SAM2 hurts Sa2VA & N/A on RES & A1: same $h$ $\to$ SAM2 is 2.43 gIoU \\
What we borrow & sparse-\Seg{} diagnosis & low-D corridor carrier & --- \\
What we do not do & linear split, CLUB, mask-reprompt & CULane/TuSimple re-implementation & --- \\
\bottomrule
\end{tabularx}
\end{table*}

\section{SAM2 zero-shot oracle: segmenter probe}
\label{app:oracle}

We additionally measured a SAM2 \citep{ravi2024sam2} zero-shot \emph{oracle}: for each test image we auto-segment an $8{\times}8$ point grid (192 candidate masks) and \emph{oracle-select} the candidate with the highest IoU against ground truth, reaching gIoU 74.88.
This oracle is \emph{not deployable}: it has no mechanism to choose the target among candidates and requires ground-truth at selection time, so we use it only as a probe of what the frozen-SAM2 segmenter can yield on this split, not as a reference row.
One reading follows: detector+SAM2 pipelines (${\approx}$25--28) leave most of this headroom unused---consistent with a bottleneck in ordinal target selection under language, not in SAM2's ability to produce a region.
We do not treat 74.88 as a reference row in the main-results table of the main paper.

\section{Literature level references (different datasets)}
\label{app:lit}

Table~\ref{tab:lit} lists public numbers from related literature; the datasets differ from FIRC2234, so these rows serve only as a \emph{level} reference, not a comparison.

\begin{table}[t]
\centering
\small
\setlength{\tabcolsep}{3pt}
\caption{Literature level references. RefCOCO / RailSem19 are different benchmarks from FIRC2234.}
\label{tab:lit}
\begin{tabularx}{\columnwidth}{@{}Ycc@{}}
\toprule
Method & Dataset & Reported metric \\
\midrule
LISA \citep{lai2024lisa} & RefCOCO val & gIoU ${\approx}74.9$ \\
GLaMM \citep{rasheed2024glamm} & RefCOCO val & gIoU ${\approx}75.0$ \\
Sa2VA-4B \citep{zhang2025sa2va} & RefCOCO & 78.9 \\
DeSa2VA-4B \citep{dang2025desa2va} & RefCOCO & 82.6 \\
DeSa2VA ablation (text pre-training w/o decoupling) & RefCOCO & 78.9$\to$73.4 \\
RailSem19-family rail segmentation \citep{zendel2019railsem19} & RailSem19 & mIoU 55--75 \\
\bottomrule
\end{tabularx}
\end{table}

\section{Label ReconIoU}
\label{app:recon}

Offline, each FIRC mask is converted to $\phih^{gt}$ (vanishing point, four centerline points, left/right widths) and rasterized with a non-learned numpy renderer used only for label QA.
Median reconstruction IoU is $0.889$ (train, $n{=}1664$), $0.892$ (train hflip), $0.890$ (val), $0.891$ (test).
The E0 gate (median ${\ge}0.85$) passes on all splits: a 20-D corridor is an adequate \emph{label} for this annotation style.
Training then uses the \emph{differentiable} $\cR$ of the rasterizer section of the main paper, which is similar in spirit (polyline + width) but is not the same renderer: $\cR$ reduces the width profile to a single scalar per side, so $\phih^{gt}$ is not the overlap-optimal parameter vector for $\cR$.
Replaying the QA labels through $\cR$ at full resolution therefore scores well below the QA figure (global IoU $0.69$ at $K{=}4$ over the $n{=}208$ test frames, against $0.891$ for the QA renderer; the renderer and the aggregation both differ, so the two figures are not directly comparable).
The ReconIoU numbers above bound the \emph{label}, not the trained decoder; the trained number is the \gIou{} in the main-results table of the main paper.

This quantization cost is deliberate and measurable rather than an artifact of the QA renderer: the 20-D carrier discards the width \emph{profile} along the corridor (the rasterizer section of the main paper), so the ${\approx}0.11$ reconstruction gap is a floor on what any decoder of $\phih$ can recover.
Continuous or higher-dimensional carriers, for instance the interpolated token space of \citet{li2026momentorpp}, trade parameters for exactly this fidelity, and we accept the trade because for a corridor target class $\phih^{gt}$ is already the annotation.

\section{Inference cost of the ECF-native path}
\label{app:cost}

Table~\ref{tab:cost-supp} reports the measured cost of the stage that replaces SAM2 mask decoding.
All rows are fp32, batch $1$, $550{\times}550$ input, on a single A100-class GPU; the ParamHead and the rasterizer are timed individually, and the end-to-end figure is a full \texttt{predict\_ecf\_mask} call after a warm-up iteration, over $3$ timed repetitions.
The rasterizer is a fixed, parameter-free function of $\phih$, so its ``parameters'' entry is the cost of the module and not of a learned tensor.

\begin{table}[t]
\centering
\small
\setlength{\tabcolsep}{4pt}
\caption{Inference cost of the ECF-native decoder, fp32, batch $1$, $550{\times}550$, one A100-class GPU. The end-to-end figure is a full decoder call and includes the shared MLLM forward pass, which dominates it and is identical across carriers; the mask-production path (parameterization, rasterization, thresholding) accounts for $2.42\%$ of it.}
\label{tab:cost-supp}
\begin{tabularx}{\columnwidth}{@{}Yrr@{}}
\toprule
Component & Parameters & Latency (ms) \\
\midrule
ParamHead $f_\psi$ ($h\!\to\!\phih$) & $4.75$\,M & $2.30$ \\
Differentiable rasterizer $\cR$ & --- (fixed) & $2.90$ \\
\midrule
ECF-native end to end & $4.75$\,M & $214.6 \pm 14.2$ \\
\quad as share of backbone & $0.205\%$ & $2.42\%$ \\
\midrule
Peak allocated VRAM & \multicolumn{2}{r}{$8.92$\,GB} \\
\bottomrule
\end{tabularx}
\end{table}

Two readings of these numbers matter for the paper's claim.
First, the mask-side trainable footprint is $4.75$\,M parameters against a $2316.2$\,M-parameter backbone ($0.205\%$), because the rasterizer is fixed: replacing the dense SAM2 mask decoder with $\cR$ removes a decoder from the trainable graph rather than adding one.
Second, the mask-production stage is $2.42\%$ of the measured end-to-end latency, so the carrier is not the runtime bottleneck: reducing the dimension of $\phih$ cannot yield meaningful savings, and the remaining cost is the shared MLLM forward pass that every carrier pays.
The LoRA-rank sweep of Table~\ref{tab:lora} shows the same for the adapters: halving the rank halves LoRA parameters at a cost of $1.7$ gIoU on the checkpoint it was measured on.

\begin{table}[t]
\centering
\small
\setlength{\tabcolsep}{4pt}
\caption{LoRA-rank sweep on an early ECF checkpoint, same fp32 $550{\times}550$ protocol. Reported to show that the trainable budget is tunable; the ranks were not re-evaluated on the final augmented checkpoint, so the \gIou{} and latency columns are not comparable to the main-results table of the main paper.}
\label{tab:lora}
\begin{tabularx}{\columnwidth}{@{}r r r r@{}}
\toprule
Rank & LoRA params & \gIou & Latency (ms) \\
\midrule
$128$ & $147.7$\,M & $49.12$ & $236.3$ \\
$64$  & $73.9$\,M  & $47.43$ & $189.7$ \\
$32$  & $36.9$\,M  & $47.03$ & $196.7$ \\
$16$  & $18.5$\,M  & $48.47$ & $204.0$ \\
\bottomrule
\end{tabularx}
\end{table}

\section{Hyperparameters}
\label{app:hparams}

\textbf{Architecture.} InternVL3-2B with the SAM2 codebase; the vision encoder is frozen.
LoRA has rank $128$, $\alpha{=}256$, dropout $0.05$, and is applied to the language model's linear projections; the embedding table and the language-model head remain trainable.
ParamHead: $1536 \to 1536 \to 1536 \to 20$, ReLU, last-layer weight zeros, identity-safe bias.
Rasterizer temperature $\tau{=}2$, densify samples $32$.
We add no extra corridor vocabulary token: an additional embedding row would be untrained against the pretrained vocabulary and would perturb tokenizer alignment.

\textbf{Loss.} $\lambda_{\mathrm{mask}}{=}1.0$, $\lambda_{\mathrm{par}}{=}0.5$, $\lambda_{\mathrm{topo}}{=}0.2$ (the topology logits are not used as a decision head; Verifier off), $\lambda_{\mathrm{ver}}{=}0.3$ with Verifier weights frozen.
A weak width-monotonicity term ($0.1$) and a far-row mask weight ($\gamma{=}1.0$) are present in the reported run and are \emph{not} claimed; a later, heavier physics-oriented recipe (gated far-weight, retuned $\gamma$) did not improve the far field and is not a main result.
When SAM2 decoding is bypassed, the SAM2 mask and dice losses are replaced by a finite zero, so that no undefined gradients propagate through the unused branch.

\textbf{Optimization.} AdamW ($\beta{=}(0.9,0.999)$), lr $4{\times}10^{-5}$, weight decay $0.05$, grad clip $0.5$, AMP fp16, batch 1, warmup 5\% + cosine to 0, 4 epochs $\times$ 10 repeats of 3328 train frames (1664 original $\cup$ 1664 hflip) $= 133{,}120$ iterations, initialized from the official InternVL3-2B Sa2VA checkpoint; seed 42.
This is a single training schedule.

\textbf{A2.} The $\lambda_{\mathrm{mask}}{=}0$ row is an architecture-identical retrain under the same orig$\cup$hflip schedule and budget, not a decode-mode switch of the main weights; the only difference from the main row is the presence of the raster mask loss.

\section{Reproducibility of the evaluation protocol}
\label{app:repro}

ECF-native numbers are generated in fp32: take the first \Seg{} hidden state, apply ParamHead $\to \cR$, threshold at $0.5$.
A1 uses the same LoRA and ParamHead weights but routes the \Seg{} embedding through SAM2, also in fp32, and detector reference rows use the same 208-frame test split, original captions, and mask-IoU protocol.
fp16 evaluation is never reported: under fp16, \Seg{} decoding collapses in our runs.

\section{Scope of the claims, limitations, and open directions}
\label{app:limits}

The main paper keeps this material out of its body; it is collected here so that the scope of each claim is on the record.

\textbf{Limitations.}
(1)~ECF assumes a corridor-shaped target; it is not a general RES decoder for compact instances.
(2)~The carrier-versus-fine-tuning separation rests on A1, which reuses our trained weights and changes only the decoder; a Sa2VA fine-tuned with its native mask loss on FIRC is not evaluated.
(3)~FIRC2234 carries a single caption template, so we can show robustness to paraphrase but not \emph{sensitivity} to a swapped target caption; a caption-swap probe needs multi-rail frames and remains the most important missing experiment.
(4)~Topology logits are reserved but unused as a fork verifier; multi-rail disambiguation is not evaluated.
(5)~Transfer sets are eval-only fill-bands with a shared caption template, not independently annotated referring corpora.
(6)~We do not re-implement DeSa2VA or parametric lane detectors on FIRC; the connection is mechanistic, not a head-to-head comparison.
(7)~A weak width-monotonicity term is included but not claimed; a later, heavier physics recipe failed to help and is not a main result.
(8)~Strict \PrAt{0.7} (87.98) sits 9.6 points below \PrAt{0.5} (97.6); residual errors concentrate on far-field thin structure.
(9)~All main results come from a single training run on a single backbone (InternVL3-2B Sa2VA weights); run-to-run variance is not reported and the detector rows use a fixed leftmost-box rule (alternatives not swept).

\textbf{Open directions.}
Five lines follow from these limits: (i)~a caption-swap probe on multi-rail frames, for linguistic sensitivity rather than paraphrase robustness; (ii)~a Rail-Path Verifier on the topology logits for forks; (iii)~DeSa2VA-style split heads so that $\phih$ comes from a visual stream, not an entangled $h$; (iv)~dropping SAM2 from the training graph entirely; and (v)~video rails, where a corridor field has a natural temporal smoothness prior.
A natural next system would combine the two treatments: DeSa2VA-style split heads so that language remains rich, plus $\cR$ so that the visual stream carries a corridor field, with a verifier for forks that the two topology logits currently ignore.
Orthogonal work can still improve the \emph{training signal}, e.g.\ token-weighted GRPO for video reasoning with focused thinking \citep{dang2025twgrpo}; the main paper leaves RL aside and changes the mask carrier.
